\documentclass[12pt,letterpaper]{article}

% ===== MCM/ICM Contest Paper Template =====
% 美赛格式规范：
%   - US Letter paper, 1-inch margins
%   - Title: 14pt bold centered
%   - Abstract/Summary: required, labeled "Summary"
%   - Section headings: 12pt bold, numbered (1, 1.1, 1.1.1)
%   - Body text: 12pt, 1.5 line spacing
%   - Tables: three-line style
%   - Figures: centered, captions below
%   - References: [1] format

% ---------- Page Setup ----------
\usepackage[top=1in,bottom=1in,left=1in,right=1in]{geometry}

% ---------- Math ----------
\usepackage{amsmath,amssymb,amsthm}

% ---------- Graphics & Tables ----------
\usepackage{graphicx,subfig}
\usepackage{booktabs,multirow,array}
\usepackage{float}

% ---------- Algorithms ----------
\usepackage{algorithm,algorithmic}
\renewcommand{\algorithmicrequire}{\textbf{Input:}}
\renewcommand{\algorithmicensure}{\textbf{Output:}}

% ---------- Hyperlinks ----------
\usepackage[colorlinks=true,linkcolor=black,citecolor=black,urlcolor=black]{hyperref}

% ---------- Lists ----------
\usepackage{enumitem}

% ---------- Captions ----------
\usepackage{caption}
\captionsetup{font=small,labelfont=bf,labelsep=quad,justification=centering}

% ---------- Headers/Footers ----------
\usepackage{fancyhdr}
\pagestyle{fancy}
\fancyhf{}
\fancyhead[L]{\small Team \#\textlangle Your Team Number\textrangle}
\fancyhead[R]{\small\thepage}
\renewcommand{\headrulewidth}{0.4pt}

% ---------- Paragraph ----------
\usepackage{indentfirst}
\setlength{\parindent}{0pt}           % MCM typically no indent
\setlength{\parskip}{6pt}              % space between paragraphs
\linespread{1.5}                       % 1.5 line spacing

% ---------- Section Format ----------
\usepackage{titlesec}
\titleformat{\section}{\bfseries\large}{\thesection}{1em}{}
\titleformat{\subsection}{\bfseries\normalsize}{\thesubsection}{1em}{}
\titleformat{\subsubsection}{\bfseries\normalsize}{\thesubsubsection}{1em}{}

% Numbering
\numberwithin{equation}{section}
\numberwithin{figure}{section}
\numberwithin{table}{section}

\begin{document}

% ==========================================
% Title
% ==========================================
\begin{center}
{\Large\bfseries\textlangle Your Paper Title Here\textrangle}

\vspace{0.3cm}
\textlangle Date\textrangle
\end{center}

\vspace{0.5cm}

% ==========================================
% Summary (Required by MCM/ICM)
% ==========================================
\noindent\textbf{\large Summary}

\vspace{0.3cm}

\textlangle Write your summary here. The summary should be a concise description of your approach, methods, and key results. It should not exceed one page. Key elements to include:
\begin{enumerate}
    \item Problem background and restatement
    \item Assumptions and justification
    \item Model(s) developed and methods used
    \item Key results and conclusions
    \item Strengths and limitations
\end{enumerate}
\textrangle

\vspace{0.5cm}
\noindent\textbf{Keywords:} keyword1; keyword2; keyword3; keyword4; keyword5

\newpage

% ==========================================
% Table of Contents
% ==========================================
\tableofcontents
\newpage

% ==========================================
% 1. Introduction
% ==========================================
\section{Introduction}

\subsection{Problem Background}
\textlangle Describe the real-world background and significance of the problem. 300--400 words.\textrangle

\subsection{Problem Restatement}
\textlangle Restate the problem in your own academic language. Clearly list the known conditions and objectives for each sub-problem. Do not copy the problem statement verbatim. 400--500 words.\textrangle

\subsection{Literature Review}
\textlangle Briefly review existing work related to this problem. Cite key references [1][2]. 200--300 words.\textrangle

% ==========================================
% 2. Assumptions and Justifications
% ==========================================
\section{Assumptions and Justifications}

\textlangle List and justify each assumption. 6--8 assumptions typical.\textrangle

\begin{enumerate}[label=\arabic*.]
    \item \textbf{Assumption 1:} \textlangle Description\textrangle. \textbf{Justification:} \textlangle Why this assumption is reasonable\textrangle.
    \item \textbf{Assumption 2:} \textlangle Description\textrangle. \textbf{Justification:} \textlangle Why reasonable\textrangle.
    \item \textbf{Assumption 3:} \textlangle Description\textrangle. \textbf{Justification:} \textlangle Why reasonable\textrangle.
\end{enumerate}

% ==========================================
% 3. Notation and Definitions
% ==========================================
\section{Notation and Definitions}

% Parameters table
\begin{table}[H]
\centering
\caption{Model Parameters}
\label{tab:parameters}
\begin{tabular}{c|l|c}
\toprule
\textbf{Symbol} & \textbf{Description} & \textbf{Unit} \\
\midrule
$N$ & Number of nodes & -- \\
\bottomrule
\end{tabular}
\end{table}

% Decision variables table
\begin{table}[H]
\centering
\caption{Decision Variables}
\label{tab:variables}
\begin{tabular}{c|l|c}
\toprule
\textbf{Symbol} & \textbf{Description} & \textbf{Type} \\
\midrule
$x_{ij}$ & Whether arc $(i,j)$ is selected & binary \\
\bottomrule
\end{tabular}
\end{table}

% ==========================================
% 4. Model Development
% ==========================================
\section{Model Development}

\subsection{Data Preprocessing}
\subsubsection{Data Overview}
\textlangle Describe the data sources, scale, and key features. 200--300 words.\textrangle

\subsubsection{Data Cleaning and Feature Engineering}
\textlangle Describe preprocessing steps. 200--300 words.\textrangle

\subsection{Model for Problem 1}

\subsubsection{Model Formulation}

\textbf{Objective Function:}

Considering \textlangle relevant factors\textrangle, the objective is to minimize \textlangle description\textrangle:

\begin{equation}\label{eq:objective}
\min Z = \sum_{i \in N} \sum_{j \in N} c_{ij} x_{ij}
\end{equation}

where $c_{ij}$ denotes \textlangle explanation\textrangle, and $x_{ij}$ is \textlangle explanation\textrangle.

\textbf{Constraints:}

\begin{enumerate}[label=(\arabic*)]
    \item \textbf{Capacity Constraint}:
    \begin{equation}\label{eq:capacity}
        \sum_{i \in N} w_i y_{ik} \leq W_k, \quad \forall k \in K
    \end{equation}
    Constraint (\ref{eq:capacity}) ensures that the total load on each vehicle does not exceed its capacity. Violating this constraint would result in infeasible loading plans.

    \item \textbf{Service Completeness}:
    \begin{equation}\label{eq:service}
        \sum_{k \in K} y_{ik} = 1, \quad \forall i \in N
    \end{equation}
    Constraint (\ref{eq:service}) ensures every customer is served exactly once.
\end{enumerate}

\subsubsection{Solution Algorithm}

\textlangle Justify algorithm choice.\textrangle

\begin{algorithm}[H]
\caption{\textlangle Algorithm Name\textrangle}
\label{alg:main}
\begin{algorithmic}[1]
\REQUIRE \textlangle Input parameters\textrangle
\ENSURE \textlangle Output\textrangle
\STATE Initialize solution
\FOR{each iteration}
    \STATE \textlangle Step description\textrangle
\ENDFOR
\RETURN Best solution found
\end{algorithmic}
\end{algorithm}

\subsubsection{Key Design Decisions}
\begin{enumerate}
    \item \textbf{Initial solution:} \textlangle Strategy description\textrangle
    \item \textbf{Neighborhood structure:} \textlangle Description\textrangle
    \item \textbf{Parameter settings:} \textlangle Key parameters and rationale\textrangle
\end{enumerate}

\subsubsection{Results for Problem 1}
\textlangle Present results with tables and figures.\textrangle

\subsection{Model for Problem 2}
\textlangle Same structure as Problem 1. 800--1000 words.\textrangle

\subsection{Model for Problem 3}
\textlangle Same structure as Problem 1. 800--1000 words.\textrangle

% ==========================================
% 5. Results and Discussion
% ==========================================
\section{Results and Discussion}

\subsection{Main Results}
\textlangle Present and interpret the main numerical results. Use tables and figures with proper references (see Figure \ref{fig:example}, Table \ref{tab:example}). 500--800 words.\textrangle

\subsection{Sensitivity Analysis}
\textlangle Analyze the sensitivity of results to key parameters. 400--500 words.\textrangle

\subsection{Model Validation}
\textlangle Describe how the model was validated. Compare with known solutions or simpler models. 200--300 words.\textrangle

% ==========================================
% 6. Strengths, Limitations, and Future Work
% ==========================================
\section{Strengths, Limitations, and Future Work}

\subsection{Strengths}
\begin{enumerate}
    \item \textlangle Strength 1 with specific example\textrangle
    \item \textlangle Strength 2 with specific example\textrangle
    \item \textlangle Strength 3\textrangle
\end{enumerate}

\subsection{Limitations}
\begin{enumerate}
    \item \textlangle Limitation 1 with impact assessment\textrangle
    \item \textlangle Limitation 2\textrangle
\end{enumerate}

\subsection{Future Work and Extensions}
\textlangle Discuss how the model could be extended or applied to other domains. 300--400 words.\textrangle

% ==========================================
% 7. Conclusions
% ==========================================
\section{Conclusions}
\textlangle Summarize the key contributions, findings, and implications. 200--300 words.\textrangle

% ==========================================
% References
% ==========================================
\newpage
\begin{thebibliography}{99}

\bibitem{ref1} Author1, Author2. Article Title[J]. \emph{Journal Name}, Year, Volume(Issue): Pages.

\bibitem{ref2} Author. Article Title[J]. \emph{Journal Name}, Year, Volume(Issue): Pages.

\bibitem{ref3} Author. \emph{Book Title}[M]. City: Publisher, Year.

\bibitem{ref4} Solomon M M. Algorithms for the Vehicle Routing and Scheduling Problems with Time Window Constraints[J]. \emph{Operations Research}, 1987, 35(2): 254--265.

\bibitem{ref5} Ropke S, Pisinger D. An Adaptive Large Neighborhood Search Heuristic for the Pickup and Delivery Problem with Time Windows[J]. \emph{Transportation Science}, 2006, 40(4): 455--472.

\end{thebibliography}

% ==========================================
% Appendix
% ==========================================
\newpage
\section*{Appendix}
\addcontentsline{toc}{section}{Appendix}

\subsection*{Appendix A: Parameter Data}
\textlangle Supplementary parameter tables.\textrangle

\subsection*{Appendix B: Key Code Segments}
{\small\ttfamily
\begin{verbatim}
# Paste key code here
\end{verbatim}
}

\end{document}
